%-------------------------------------------------------------------------------
% Mẫu CV / Resume Chuẩn Đại học Harvard
% Ứng viên: HUỲNH CÔNG LUYỆN
% Vị trí mục tiêu: Kỹ sư Backend AI | Kỹ sư AI Backend | Kỹ sư Backend (Hệ thống AI)
%-------------------------------------------------------------------------------

\documentclass[10pt,letterpaper]{article}

% --- Nhận diện Engine & Hỗ trợ Tiếng Việt / Unicode ---
\usepackage{iftex}
\ifPDFTeX
  \usepackage[utf8]{vietnam}
\else
  \usepackage{fontspec}
\fi

% --- Căn chỉnh lề trang (Chuẩn ATS & Khớp đúng 1 trang A4/Letter) ---
\usepackage[
  letterpaper,
  top=0.35in,
  bottom=0.35in,
  left=0.42in,
  right=0.42in
]{geometry}

% --- Gói định dạng văn bản ---
\usepackage{titlesec}
\usepackage{enumitem}
\usepackage{xcolor}
\usepackage{hyperref}
\usepackage{tabularx}
\usepackage[verbose=silent]{microtype}
\linespread{1.04}

% --- Bảng màu ---
\definecolor{primary}{RGB}{0, 0, 0}
\definecolor{secondary}{RGB}{45, 45, 45}
\definecolor{linkcolor}{RGB}{18, 65, 145}

% --- Cấu hình Hyperref ---
\hypersetup{
  colorlinks=true,
  linkcolor=linkcolor,
  filecolor=linkcolor,
  urlcolor=linkcolor,
  pdfauthor={Huỳnh Công Luyện},
  pdftitle={Huỳnh Công Luyện - CV Kỹ sư Backend AI},
  pdfsubject={CV Kỹ sư Backend AI},
  pdfkeywords={Backend, AI Systems, RAG, REST APIs, Microservices, Node.js, Spring Boot, FastAPI, Docker}
}

% --- Bỏ thụt lề đầu dòng & Số trang ---
\pagestyle{empty}
\raggedbottom
\setlength{\tabcolsep}{0in}

% --- Định dạng tiêu đề Section ---
\titleformat{\section}{
  \raggedright\large\bfseries\color{primary}
}{}{0em}{}[\vspace{-2.5pt}\color{black}\titlerule \vspace{2pt}]
\titlespacing*{\section}{0pt}{7.5pt plus 1pt minus 1pt}{3.5pt plus 1pt minus 1pt}

% --- Định dạng danh sách ---
\newcommand{\resumeItemListStart}{\begin{itemize}[leftmargin=1.15em, label={\small$\bullet$}, itemsep=1.2pt, parsep=0pt, topsep=1.0pt, partopsep=0pt]}
\newcommand{\resumeItemListEnd}{\end{itemize}\vspace{-1pt}}

\newcommand{\resumeItem}[1]{\item \small{#1}}

\newcommand{\resumeSubheading}[4]{
  \item[]
  \begin{tabular*}{\textwidth}[t]{l@{\extracolsep{\fill}}r}
    \textbf{#1} & \textbf{\small #2} \\
    \textit{\small #3} & \textit{\small #4} \\
  \end{tabular*}\vspace{0.5pt}
}

\begin{document}

%===============================================================================
% THÔNG TIN CÁ NHÂN / LIÊN HỆ
%===============================================================================
{\centering
  {\Large \textbf{Huỳnh Công Luyện}}\\[2pt]
  {\large \textbf{\color{secondary}KỸ SƯ BACKEND AI}} \quad{\small$|$\quad \textbf{Kiến trúc Backend $\bullet$ REST APIs $\bullet$ Tích hợp AI/LLM $\bullet$ Hệ thống Truy xuất}}\\[3pt]
  \small{
    Email: \href{mailto:huynhluyen839@gmail.com}{huynhluyen839@gmail.com} \quad$\bullet$\quad
    Điện thoại: (+84) 926072506 \quad$\bullet$\quad
    Địa điểm: TP. Hồ Chí Minh, Việt Nam
  }\\[1.8pt]
  \small{
    GitHub: \href{https://github.com/neshaki091}{github.com/neshaki091} \quad$\bullet$\quad
    Portfolio: \href{https://Neshaki091.vercel.app}{Neshaki091.vercel.app} \quad$\bullet$\quad
    Trạng thái: Sẵn sàng làm việc toàn thời gian
  }\par
}
\vspace{2pt}

%===============================================================================
% TÓM TẮT CHUYÊN MÔN
%===============================================================================
\section{Tóm tắt Chuyên môn}
\small{
Ứng viên IT (Đã hoàn thành chương trình, GPA: 3.34/4.0) có kinh nghiệm thực tế về thiết kế dịch vụ backend, tích hợp LLM API và xây dựng ứng dụng RAG end-to-end. Tác giả nền tảng \textbf{Waterbase} (hệ thống BaaS đa người dùng tự host) và \textbf{ARQ-RAG} (engine truy xuất vector tối ưu hóa bộ nhớ). Thành thạo Node.js, Python/FastAPI, Java/Spring Boot, cơ sở dữ liệu, Docker và triển khai cloud.
}

%===============================================================================
% HỌC VẤN
%===============================================================================
\section{Học vấn}
\begin{itemize}[leftmargin=0pt, label={}, topsep=1.0pt, partopsep=0pt, parsep=0pt]
  \resumeSubheading
    {Trường Đại học Giao thông vận tải TP. Hồ Chí Minh (UTH)}{TP. Hồ Chí Minh, Việt Nam}
    {Cử nhân Công nghệ Thông tin \textnormal{(Đã hoàn thành chương trình; Chờ cấp bằng)}}{2022 -- 2026}
    \resumeItemListStart
      \resumeItem{\textbf{GPA:} 3.34 / 4.0 \quad$\bullet$\quad \textbf{Khóa luận:} ARQ-RAG --- Điểm bảo vệ: 9.0 / 10.0}
    \resumeItemListEnd
\end{itemize}

%===============================================================================
% KỸ NĂNG CỐT LÕI
%===============================================================================
\section{Kỹ năng Cốt lõi}
\begin{itemize}[leftmargin=0pt, label={}, itemsep=1.2pt, topsep=1.0pt, partopsep=0pt, parsep=0pt]
  \item \small{\textbf{Phát triển Backend:} Node.js, Express, Python, FastAPI, Java, Spring Boot, REST APIs, Xử lý I/O bất đồng bộ}
  \item \small{\textbf{Cơ sở Dữ liệu \& Hệ thống:} MongoDB, Redis, Qdrant, Supabase pgvector, chỉ mục vector, zero-copy \texttt{mmap}}
  \item \small{\textbf{Hạ tầng \& Bảo mật:} Docker Compose, Nginx, JWT, Refresh-token rotation, Phân quyền RBAC, GCP Cloud Functions, Git}
  \item \small{\textbf{Hệ thống AI:} Tích hợp API LLM, Pipeline RAG, Embeddings, Truy xuất ngữ nghĩa, Prompt Engineering, Đánh giá RAG, QLoRA}
  \item \small{\textbf{Ngoại ngữ:} Tiếng Anh: A2 (đọc hiểu tài liệu kỹ thuật và papers) $\bullet$ Tiếng Việt: Bản ngữ}
\end{itemize}

%===============================================================================
% DỰ ÁN & NGHIÊN CỨU TIÊU BIỂU
%===============================================================================
\section{Dự án \& Nghiên cứu Tiêu biểu}
\begin{itemize}[leftmargin=0pt, label={}, topsep=1.0pt, partopsep=0pt, parsep=0pt]

  % --- DỰ ÁN 1: EduChatbot ---
  \item[]
  \begin{tabular*}{\textwidth}[t]{l@{\extracolsep{\fill}}r}
    \textbf{\large EduChatbot --- Trợ lý Hỏi đáp Tài liệu RAG Tiếng Việt} & \textbf{\small Thg 06/2026} \\
    \textit{\small Trưởng nhóm (6 thành viên) \quad$|$\quad Java 17, Spring Boot, Python, FastAPI, Qdrant, Llama 3.1 8B \quad$|$\quad \href{https://github.com/Neshaki091/Educhatbot}{GitHub}} & \textit{\small Đồ án Môn học Java} \\
  \end{tabular*}\vspace{0.5pt}
  \resumeItemListStart
    \resumeItem{\textbf{Kiến trúc Backend \& Tích hợp:} Xây dựng các service core Java 17 / Spring Boot với Spring Security JWT và phân quyền chi tiết (\texttt{@PreAuthorize}); tích hợp cầu nối bất đồng bộ Python/FastAPI qua Spring WebClient để phân tách mô hình LLM và xử lý inference.}
    \resumeItem{\textbf{Tích hợp AI, Đánh giá \& Triển khai:} Xử lý pipeline nạp tài liệu tiếng Việt, embedding 384 chiều và truy xuất vector Qdrant; trực tiếp deploy backend lên GCP Cloud Functions với Supabase + Qdrant Cloud; điều phối nhóm 6 người và thực hiện đánh giá mù A/B Ragas (faithfulness 0,303 vs. 0,292 so với fine-tuning).}
  \resumeItemListEnd
  \vspace{4.0pt}

  % --- DỰ ÁN 2: Waterbase ---
  \item[]
  \begin{tabular*}{\textwidth}[t]{l@{\extracolsep{\fill}}r}
    \textbf{\large Waterbase --- Nền tảng BaaS Đa người dùng Tự host} & \textbf{\small Thg 11/2025} \\
    \textit{\small Lập trình viên duy nhất \quad$|$\quad Thiết kế Kiến trúc, Lập trình \& Triển khai \quad$|$\quad \href{https://tttn-project-waterbase-mini-clone-f.vercel.app/}{Live Demo}} & \textit{\small Đồ án Thực tập Tốt nghiệp} \\
  \end{tabular*}\vspace{0.5pt}
  \resumeItemListStart
    \resumeItem{\textbf{Kiến trúc Backend \& Triển khai:} Thiết kế và triển khai nền tảng microservices gồm 7 services (Auth, App, WaterDB, Realtime, Storage, Rule Engine, Analytics) bằng Node.js, Express, MongoDB, Redis và Docker Compose trên VPS tối thiểu 1 GB RAM / 1 vCPU.}
    \resumeItem{\textbf{Multi-Tenancy \& Bảo mật:} Hiện thực cô lập dữ liệu theo mô hình Database-per-Application triệt tiêu rủi ro rò rỉ dữ liệu, xác thực JWT kèm xoay vòng refresh token, và Rule Engine (RBAC) phân quyền khai báo chi tiết cấp collection.}
    \resumeItem{\textbf{Hạ tầng Realtime \& API:} Xây dựng đồng bộ dữ liệu thời gian thực qua WebSocket (RTWaterDB) và cấu hình Nginx làm API Gateway cho reverse proxy, SSL termination và định tuyến microservice.}
  \resumeItemListEnd
  \vspace{4.0pt}

  % --- DỰ ÁN 3: ARQ-RAG ---
  \item[]
  \begin{tabular*}{\textwidth}[t]{l@{\extracolsep{\fill}}r}
    \textbf{\large ARQ-RAG --- Engine Truy xuất Vector Thích ứng} & \textbf{\small 2026} \\
    \textit{\small Tác giả chính \quad$|$\quad Hạ tầng Truy xuất \& Hiện thực Rust bằng AI \quad$|$\quad \href{https://arq-rag-frontend.vercel.app/}{Live Demo}} & \textit{\small Khóa luận: \textbf{9.0 / 10.0}} \\
  \end{tabular*}\vspace{0.5pt}
  \resumeItemListStart
    \resumeItem{\textbf{Kiến trúc Hệ thống:} Thiết kế chỉ mục IVF, lượng tử hóa phần dư và định tuyến cụm thích ứng; hiện thực và xác thực engine truy xuất Rust với SIMD AVX2 và memory-mapped I/O (\texttt{mmap}) thông qua hỗ trợ từ AI coding agent.}
    \resumeItem{\textbf{Tối ưu Bộ nhớ:} Giảm RAM hoạt động từ 1.931 MB (FAISS baseline) xuống \textbf{102,6 MB} ($18,8\times$) trên 5 triệu vector Wikipedia DPR (768D, CPU đơn luồng) bằng cách chỉ giữ mã 4-bit trên RAM, phục vụ vector gốc từ đĩa qua zero-copy \texttt{mmap}.}
    \resumeItem{\textbf{Hiệu năng:} Đạt \textbf{4.522 QPS} tại 87,2\% Recall@16 trước rerank và \textbf{99,2\% Recall@16} sau khi Float32 rerank tại 3.450 QPS. Lập chỉ mục 500K vector trong 20,2s (nhanh hơn $1,93\times$ so với FAISS IVFPQ).}
    \resumeItem{\textbf{Nghiên cứu:} Tác giả bài báo nghiên cứu đã submit tới hội nghị AMI 2026 [Đang bình duyệt]; thiết lập pipeline benchmark có khả năng tái tạo và kỹ thuật gom cụm thân thiện với CPU cache.}
  \resumeItemListEnd
  \vspace{4.0pt}

  % --- DỰ ÁN 4: Trợ lý AI Nội bộ ---
  \item[]
  \begin{tabular*}{\textwidth}[t]{l@{\extracolsep{\fill}}r}
    \textbf{\large Trợ lý AI Nội bộ --- Lotus Hackathon VNG 2026} & \textbf{\small Thg 03/2026} \\
    \textit{\small Lập trình viên Backend (đội 2 người) \quad$|$\quad Node.js, Express, Supabase pgvector, Gemini API, n8n \quad$|$\quad \href{https://devpost.com/software/internal-ai-assistance}{Video Demo}} & \textit{\small Tập đoàn VNG} \\
  \end{tabular*}\vspace{0.5pt}
  \resumeItemListStart
    \resumeItem{\textbf{Tích hợp AI Backend:} Xây dựng backend Node.js/Express tích hợp Gemini API, Supabase \texttt{pgvector} và tự động hóa n8n, cô lập đa khách hàng theo \texttt{clientId} và router dự phòng ý định 3 tầng (parser $\rightarrow$ Gemini $\rightarrow$ từ khóa); bàn giao demo hoạt động trong \textbf{36 giờ}.}
  \resumeItemListEnd

\end{itemize}

\end{document}
